\documentclass[pdf,slideColor,fyma]{prosper}
\usepackage[czech]{babel}
\usepackage[utf8]{inputenc}
\usepackage{color}
\usepackage{graphics}
\usepackage{listings}

\lstloadlanguages{Ruby}
\lstset{
basicstyle=\ttfamily\color{black},
commentstyle = \ttfamily\color{red},
keywordstyle=\ttfamily\color{blue},
stringstyle=\color{orange}}

\slideCaption{\textit{Ruby\dots}}
\DefaultTransition{Wipe}

\title{\Large Ruby}
\author{Daniel Klimaj}
\date{2015}
\email{xklima22@stud.fit.vutbr.cz}
\institution{Vysoké učení technické, Fakulta informačních technologií}

\begin{document}
\maketitle

\begin{slide}{Ruby}
\begin{itemize}
  \item{Prvé vydanie v roku 1995}
\end{itemize}
\begin{itemize}
  \item{Autor  Yukihiro "Matz" Matsumoto}
\end{itemize}
\begin{itemize}
  \item{Inšpirovaný jazykmi C, Perl, Smalltalk, Eiffel, Ada a Lisp}
\end{itemize}
\begin{center}
  \includegraphics[width=3cm]{ruby.eps}
\end{center}
\end{slide}

\begin{slide}{História}
\begin{itemize}
  \item{1995: Prvé vydanie Ruby~0.95}
  \item{1996: Vychádza Ruby~1.0}
  \item{1997: Publikovaný prvý článok o Ruby}
  \item{1998: Vychádza Ruby~1.2}
  \begin{itemize}
    \item{Anglická webová stránka}
    \item{Ruby Application Archive}
    \begin{itemize}
      \item{hosting Ruby aplikácií}
      \item{ukončený v roku 2013 pre nedostatok zdrojov pre hostovanie}
    \end{itemize}
  \end{itemize}
  \item{1999: Vychádza Ruby~1.4}
  \begin{itemize}
    \item{Prvá kniha The Object-Oriented Scripting Language Ruby}
  \end{itemize}
  \item{2000: Vychádza Ruby~1.6}
  \begin{itemize}
    \item{Prvá kniha v angličtine Programming Ruby}
  \end{itemize}
\end{itemize}
\end{slide}

\begin{slide}{História}
\begin{itemize}
  \item{2003: Vychádza Ruby~1.8}
  \begin{itemize}
    \item{ISO štandard}
  \end{itemize}
  \item{2007: Vychádza Ruby~1.9}
  \begin{itemize}
    \item{Ruby~1.9 málo kompatibilný s Ruby~1.8}
    \item{Pomalšie rozšírenie kvôli portovaniu}
  \end{itemize}
  \item{2011: Vychádza Ruby~1.9.3}
  \item{2013: Vychádza Ruby~2.0}
  \begin{itemize}
    \item{Takmer plne kompatibilný s Ruby~1.9}
  \end{itemize}
  \item{2014: Vychádza Ruby~2.2}
  \item{2015: Ukončenie údržby Ruby~1.9}
\end{itemize}
\end{slide}

\begin{slide}{Sémantika}
\begin{itemize}
  \item{Každá hodnota je objekt, vrátane typov v iných jazykoch označených ako primitíva - integer, null, atď.}
  \item{Každá funkcia je metóda - je volaná nad objektom}
  \item{Dynamicky typovaný}
  \item{Premenné vždy uchovávajú odkazy (referencie)}
  \item{Dedičnosť - dynamický výber, mixin, singleton}
  \item{Viacnásobná dedičnosť nepodporaovaná}
  \item{Paradigma - procedurálny, funkcionálny, objektovo-orientovaný, imperatívny}
\end{itemize}
\end{slide}

\begin{slide}{Syntax}
\begin{itemize}
  \item{Podobná Perlu a Pythonu}
  \item{Bloky definované párom kľúčových slov begin a end, v niektorých prípadoch párom \{ a \}}
  \item{Príkazy sa ukončujú novým riadkom alebo s ;}
  \item{Všetky instančné premenné su privátne - prístup cez accessory \textit{attr\_reader}, \textit{attr\_writer} alebo \textit{attr\_accessor}}
  \item{Umožňuje na konci názvov metód používať znaky ? a !, špecialny význam}
    \begin{itemize}
        \item{? - návratová hodnota je logická hodnota}
          \begin{itemize}
            \item{napr. def equal?(other)}
          \end{itemize}
        \item{! - modifikuje volajúci objekt}
          \begin{itemize}
            \item{$x = "abc"$}
            \item{$x.sub("b", "d") \# x == "abc"$}
            \item{$x.sub!("b", "d") \# x == "adc"$}
          \end{itemize}
     \end{itemize}
\end{itemize}
\end{slide}

\overlays{3}{
\begin{slide}{Accessory}
\begin{figure}[ht]
\begin{picture}(320,200)
  \linethickness{1pt}
  \fromSlide{1}{
    \put(0,140){\framebox(120,60){
     \shortstack[l]{\textit{attr\_reader :var} \\ \\ def var \\ \hspace{0.2cm} @var \\ end}}}}
  \fromSlide{2}{
    \put(0,20){\framebox(120,60){
      \shortstack[l]{\textit{attr\_writer :var} \\ \\ def var=(value) \\ \hspace{0.2cm} @var = value \\ end}}}}
   \fromSlide{3}{
    \put(200,50){\framebox(120,120){
       \shortstack[l]{\textit{attr\_accessor :var} \\ \\ def var \\ \hspace{0.2cm} @var \\ end \\ \\ def var=(value) \\ \hspace{0.2cm} @var = value \\ end}}}}
    \end{picture}
  \end{figure}
\end{slide}
}


\begin{slide}{Ukážka kódu}

\begin{lstlisting}[basicstyle=\tiny, language=Ruby]
class ExampleBase
  def say_hello_to name
    puts "Hello #{name}!"
  end
end

class Example < ExampleBase
  attr_accessor :var

  def initialize val
    @var = val
  end
end

e = Example.new(5)
puts e.var              # 5
e.say_hello_to "World"  # Hello World!
\end{lstlisting}

\end{slide}

\begin{slide}{Zdroje}

\begin{itemize}
\item{Wikipedia \\
  \tiny\url{ http://en.wikipedia.org/wiki/Ruby_(programming_language)}}
\item{Ruby Lang \\
  \tiny\url{ https://www.ruby-lang.org/en/about/}}
\end{itemize}
\end{slide}
\end{document}